\exercisesfor{6}

\begin{enumerate}
\def\labelenumi{\arabic{enumi}.}
\tightlist
\item
  证明 \(\exp (\mathrm{U}).\exp (\mathrm{V}).\exp (-\mathrm{U}) = \exp \left(\sum_{n = 0}^{\infty}\frac{1}{n!} (\operatorname {ad}\mathrm{U})^{n}(\mathrm{V})\right)\)。推出
\end{enumerate}

\[
\mathrm{H} (\mathrm{U}, \mathrm{H} (\mathrm{V}, - \mathrm{U})) = \sum_ {n = 0} ^ {\infty} \frac {1}{n !} (\operatorname{ad} \mathrm{U}) ^ {n} (\mathrm{V}) = (\exp (\operatorname{ad} \mathrm{U})) (\mathrm{V}).
\]

\begin{enumerate}
\def\labelenumi{\arabic{enumi}.}
\setcounter{enumi}{1}
\item
  证明 \(\mathbf{H}(\mathbf{U},\mathbf{V}) = \sum_{n = 0}^{\infty}\frac{1}{n!} (\operatorname {ad}\mathbf{U})^{n}(\mathbf{H}(\mathbf{V},\mathbf{U}))\)。（应用习题 1，注意 \(\mathbf{H}(\mathbf{U},\mathbf{H}(\mathbf{H}(\mathbf{V},\mathbf{U}), - \mathbf{U})) = \mathbf{H}(\mathbf{U},\mathbf{V}).\)）
\item
  令 \(\mathbf{M}(\mathbf{U},\mathbf{V}) = \mathbf{H}(-\mathbf{U},\mathbf{U} + \mathbf{V})\)。于是
\end{enumerate}

\[
\exp (\mathrm{M} (\mathrm{U}, \mathrm{V})) = \exp (- \mathrm{U}). \exp (\mathrm{U} + \mathrm{V}).
\]

令 \(M_{1}(U, V)\)（相应地，\(H_{1}(U, V)\)）表示级数 M（相应地，H）中关于 V 的次数为 1 的双齐次项之和。

\begin{enumerate}
\def\labelenumi{(\alph{enumi})}
\tightlist
\item
  证明
\end{enumerate}

\[
\mathrm{M} _ {1} (\mathrm{U}, \mathrm{V}) = \sum_ {n = 0} ^ {\infty} \frac {(- 1) ^ {n}}{(n + 1) !} (\mathrm{adU}) ^ {n} (\mathrm{V}) = (f (\mathrm{adU})) (\mathrm{V})
\]

其中 \(f(\mathrm{T}) = (1 - e^{-\mathrm{T}}) / \mathrm{T}\)。

（论证与命题 5 的证明相同。也可以利用恒等式将其归结为命题 5：

\[
\exp (\mathrm{U}). \exp (\mathrm{M} (\mathrm{U}, \mathrm{V})). \exp (- \mathrm{U}) = \exp (\mathrm{H} (\mathrm{U} + \mathrm{V}, - \mathrm{U})),
\]

并结合习题 1。）

\begin{enumerate}
\def\labelenumi{(\alph{enumi})}
\setcounter{enumi}{1}
\tightlist
\item
  证明 \(\mathrm{H}(\mathrm{U},\mathrm{M}(\mathrm{U},\mathrm{V})) = \mathrm{U} + \mathrm{V}\)。推出
\end{enumerate}

\[
\mathrm{H} _ {1} (\mathrm{U}, \mathrm{M} _ {1} (\mathrm{U}, \mathrm{V})) = \mathrm{V}.
\]

\begin{enumerate}
\def\labelenumi{(\alph{enumi})}
\setcounter{enumi}{2}
\tightlist
\item
  令 \(g(\mathrm{T}) = 1 / f(\mathrm{T}) = 1 + \mathrm{T} / 2 + \sum_{n=1}^{\infty} \frac{1}{(2n)!} b_{2n}\mathrm{T}^{2n}\)，其中 \(b_{2n}\) 是
\end{enumerate}

伯努利数（《实变函数》，第六章，§ 1，第 4 号）。证明

\[
\mathrm{H} _ {1} (\mathrm{U}, \mathrm{V}) = (g (\operatorname{ad} \mathrm{U})) (\mathrm{V}) = \mathrm{V} + \frac {1}{2} [ \mathrm{U}, \mathrm{V} ] + \sum_ {n = 1} ^ {\infty} \frac {1}{(2 n) !} b _ {2 n} (\operatorname{ad} \mathrm{U}) ^ {2 n} (\mathrm{V}).
\]

写出 \(H_{1}(U, V)\) 展开的前几项：

\[
\begin{array}{r l} \mathrm{H} _ {1} (\mathrm{U}, \mathrm{V}) & = \mathrm{V} + \frac {1}{2} (\mathrm{adU}) (\mathrm{V}) + \frac {1}{12} (\mathrm{adU}) ^ {2} (\mathrm{V}) - \frac {1}{720} (\mathrm{adU}) ^ {4} (\mathrm{V}) \\ & + \frac {1}{30240} (\mathrm{adU}) ^ {6} (\mathrm{V}) - \dots \end{array}
\]

\begin{enumerate}
\def\labelenumi{(\alph{enumi})}
\setcounter{enumi}{3}
\tightlist
\item
  证明 H(U, V) 中关于 U 的次数为 1 的双齐次项之和为级数：
\end{enumerate}

\[
\mathrm{U} + \frac {1}{2} [ \mathrm{U}, \mathrm{V} ] + \sum_ {n = 1} ^ {\infty} \frac {1}{(2 n) !} b _ {2 n} (\mathrm{adV}) ^ {2 n} (\mathrm{U}).
\]

（将 (c) 应用于 H(V, U)，并使用习题 2 将 H(V, U) 转换为 H(U, V)。）(e) 令 W 为另一个不定元，令 X(U, V, W) 为 H(U, V + W) 中关于 W 的次数为 1 的三齐次项之和。证明

\[
\mathrm{X} (\mathrm{U}, \mathrm{H} _ {1} (\mathrm{V}, \mathrm{W})) = \mathrm{H} _ {1} (\mathrm{H} (\mathrm{U}, \mathrm{V}), \mathrm{W}).
\]

（使用恒等式 \(\mathrm{H}(\mathrm{U},\mathrm{H}(\mathrm{V},\mathrm{W})) = \mathrm{H}(\mathrm{H}(\mathrm{U},\mathrm{V}),\mathrm{W}).\)）推出

\[
\mathrm{X} (\mathrm{U}, \mathrm{V}, \mathrm{W}) = (g (\operatorname{ad} \mathrm{H} (\mathrm{U}, \mathrm{V})) \circ f (\operatorname{ad} \mathrm{V})) (\mathrm{W}).
\]

¶ 4. 令 X 是有限集，令 \(\hat{\mathbf{L}} = \hat{\mathbf{L}}(\mathbf{X})\)。在 \(\hat{\mathbf{L}}\) 上赋予 Hausdorff 复合律 \((a, b) \mapsto a \uplus b\)，并简记为 \(a.b\)。若 \(t \in \mathbf{K}\)、\(a \in \hat{\mathbf{L}}\)，则记 \(a^t = t a\)。

\begin{enumerate}
\def\labelenumi{(\alph{enumi})}
\tightlist
\item
  令 \(\mathcal{H}\) 是关于 X 的 Hall 集。若 \(m \in \mathcal{H}\)，令 \(m_{\mathrm{L}}\) 表示李代数 \(\mathrm{L}(\mathrm{X})\) 中相应的元素（§2，第 11 号），令 \(m_{\mathcal{H}}\) 表示群 \(\hat{\mathrm{L}}\) 的基本交换子，它由 \(m\) 定义（§5，第 4 号，注）。若 \(m\) 的长度为 \(l\)，证明 \(m_{\mathcal{H}} = m_{\mathrm{L}} + \mu\)，其中 \(\mu\) 的滤过次数 \(\geqslant l + 1\)（对 \(l\) 作归纳）。推出每个元素 \(w\) 都属于群 \(\hat{\mathrm{L}}\)，并且能唯一地写成
\end{enumerate}

\[
w = \prod_ {m \in \mathcal {H}} m _ {\mathcal {H}} ^ {\alpha (m)},
\]


其中 \(\alpha(m) \in \mathbf{K}\)，且该乘积在拓扑群 \(\hat{\mathbf{L}}\) 中收敛。

\begin{enumerate}
\def\labelenumi{(\alph{enumi})}
\setcounter{enumi}{1}
\item
  令 \(\mathbf{P}\) 为素数集合，令 \(c\) 为满足 \(\geqslant 1\) 的整数。令 \(\mathbf{K} = \mathbf{Q}\)。证明群 \(\hat{\mathbf{L}} / \hat{\mathbf{L}}_c\) 可识别为 \(\mathbf{P}\)-包络 \(\mathrm{F}_{\mathbf{P}}^{c}\)，其所包络的群为 \(\mathrm{F}^c = \mathrm{F}(\mathrm{X}) / \mathrm{C}^c\mathrm{F}(\mathrm{X})\)，参见~§ 5，习题 6。
\item
  取 X 为含两个元素 \(\{U, V\}\) 的集合。证明存在两个由有理数组成的族 \(\alpha(m)\)、\(\beta(m)\)，使得
\end{enumerate}

\[
\mathrm{U} + \mathrm{V} = \prod_ {m \in \mathcal {H}} m _ {\mathcal {H}} ^ {\alpha (m)}
\]

\[
[ \mathrm{U}, \mathrm{V} ] = \prod_ {m \in \mathcal {H}} m _ {\mathcal {H}} ^ {\beta (m)}.
\]

例如

\[
\begin{array}{c} \mathrm{U} + \mathrm{V} = \mathrm{U}. \mathrm{V}. (\mathrm{U}, \mathrm{V}) ^ {- \frac {1}{2}} \dots \\ {} [ \mathrm{U}, \mathrm{V} ] = (\mathrm{U}, \mathrm{V}) \dots \end{array}
\]

\begin{enumerate}
\def\labelenumi{(\alph{enumi})}
\setcounter{enumi}{3}
\tightlist
\item
  令 g 是带有 Hausdorff 群律的幂零李 Q-代数，令 u、v 属于 g。证明
\end{enumerate}

\[
u + v = \prod_ {m \in \mathcal {H}} m _ {\mathcal {H}} (u, v) ^ {\alpha (m)}
\]

\[
[ u, v ] = \prod_ {m \in \mathcal {H}} m _ {\mathcal {H}} (u, v) ^ {\beta (m)},
\]

其中 \(\alpha\) 和 \(\beta\) 是上面定义的有理数族（“Hausdorff 公式的逆”）。（使用连续同态 \(\phi\colon\hat{L}\to\mathfrak{g}\)，满足 \(\phi(U)=u\)、\(\phi(V)=v\)，并注意到它对于 Hausdorff 群律也是同态。）

\begin{enumerate}
\def\labelenumi{(\alph{enumi})}
\setcounter{enumi}{4}
\tightlist
\item
  令 G 是无挠可除幂零群（即 P-无挠且 P-可除）。若 \(u \in G\)、\(t \in Q\)，则 \(u^{t}\) 已定义（\(\S 4\)，习题 16）。若 \(u \in G\)、\(v \in G\)，则用上面的 (d) 中的公式定义 \(u + v\) 和 {[}u, v{]}。证明这样 G 具有幂零李 Q-代数结构，且相应的 Hausdorff 律就是 G 上给定的群律（当 G 是群 \(F_{P}^{c}\) 时验证这些断言，参见~(b)），再利用从 \(F_{P}^{r}\) 到 G 的同态将其推广到一般情形。
\end{enumerate}

令 \(f\) 是从 \(G\) 到无挠可除幂零群 \(G'\) 的映射。证明 \(f\) 是（群）同态，当且仅当它是李代数同态。（因此，Hausdorff 律在幂零李 \(\mathbf{Q}\)-代数范畴与无挠可除幂零群范畴之间定义了一个同构。）

\begin{enumerate}
\def\labelenumi{\arabic{enumi}.}
\setcounter{enumi}{4}
\tightlist
\item
  令 \(\mathfrak{g}\) 是赋予 Hausdorff 群律的幂零李 \(\mathbf{Q}\)-代数，将该群律记为 \((x,y)\mapsto x.y\)。
\end{enumerate}

\begin{enumerate}
\def\labelenumi{(\alph{enumi})}
\tightlist
\item
  证明若 \(x \in \mathfrak{g}\)、\(y \in \mathfrak{g}\)，则
\end{enumerate}

\[
x. y. x ^ {- 1} = \sum_ {n = 0} ^ {\infty} \frac {1}{n !} (\operatorname{ad} x) ^ {n} (y).
\]

（使用习题 2。）

\begin{enumerate}
\def\labelenumi{(\alph{enumi})}
\setcounter{enumi}{1}
\item
  令 \(\mathfrak{h}\) 是 \(\mathfrak{g}\) 的子集。证明 \(\mathfrak{h}\) 是 \(\mathfrak{g}\) 的李子代数（相应地，理想），当且仅当它是群 \(\mathfrak{g}\) 的饱和子群（相应地，饱和正规子群）。（使用习题 4(d) 的公式，将群律转换为李代数律。）
\item
  令 \(\mathfrak{h}\) 是群 \(\mathfrak{g}\) 的子群。证明 \(\mathbf{P}\)-饱和的 \(\mathfrak{h}\) 是 \(\mathbf{Q}.\mathfrak{h}\)。
\item
  令 \(\mathfrak{h}\) 是 \(\mathfrak{g}\) 的李子代数。证明 \(\mathfrak{h}\) 在群 \(\mathfrak{g}\) 中的中心化子（相应地，正规化子）是所有 \(x \in \mathfrak{g}\) 中满足 \((\operatorname{ad} x)(\mathfrak{h}) = 0\)（相应地，\((\operatorname{ad} x)(\mathfrak{h}) \subset \mathfrak{h}\)）者所成的集合。
\item
  证明李代数 g 的下中心列与群 g 的下中心列重合，并且相关的分次李代数 gr(g) 从“群”的观点看与从“李代数”的观点看相同。
\end{enumerate}

¶ 6. 令 G 是有限生成的无挠幂零群。证明当 n 充分大时，G 可嵌入 Z 上的 n 阶严格下三角群中。（令 \((i, \overline{G})\) 为 G 的一个 P-包络，参见~§ 5，习题 6；赋予其典范李代数结构后（习题 4），\(\overline{G}\) 是有限维幂零李 Q-代数。将 Ado 定理（第一章，§ 7）应用于 \(\overline{G}\)，推出 \(\overline{G}\) 可嵌入 Q 上的严格下三角群。再通过适当的整对角矩阵共轭，将其从 Q 转到 Z。）

